\documentclass[11pt,reqno]{amsart}

\usepackage{amsmath,amssymb,amsthm}
\usepackage[margin=1.15in]{geometry}
\usepackage{microtype}

\newtheorem{theorem}{Theorem}
\newtheorem{proposition}[theorem]{Proposition}

\newcommand{\N}{\mathbb N}
\newcommand{\Z}{\mathbb Z}
\newcommand{\card}[1]{\lvert #1\rvert}

\title{The Erd\H{o}s--Graham irregularity problem for subset sums}
\author{Principia Math}
\date{}

\begin{document}

\begin{abstract}
For $c>0$, let $f_c(n)$ be the largest cardinality of a set
$A\subseteq\{1,\ldots,\lfloor cn\rfloor\}$ for which $n$ is not a subset sum of $A$.
We prove that $f_c(n)/n$ does not converge for every $0<c<1$.  For $c\geq1$ we have the exact formula
$f_c(n)=\lfloor cn\rfloor-\lceil n/2\rceil$, and hence $f_c(n)/n\to c-1/2$.
The only non-elementary input is a bounded zero-sum theorem of Alon.
\end{abstract}

\maketitle

For a finite set $A$ of positive integers, write
\[
 \Sigma(A)=\left\{\sum_{a\in B}a:B\subseteq A\right\},
 \qquad
 F(M,n)=\max\{\card A:A\subseteq[1,M],\ n\notin\Sigma(A)\}.
\]
Thus $f_c(n)=F(\lfloor cn\rfloor,n)$.  Erd\H{o}s and Graham asked whether this extremal quantity depends irregularly on $n$ \cite[p.~59]{EG}.  We prove the following complete dichotomy.

\begin{theorem}\label{thm:main}
Let $c>0$.
\begin{enumerate}
\item If $0<c<1$, then $f_c(n)/n$ does not converge.
\item If $c\geq1$, then
\[
 f_c(n)=\lfloor cn\rfloor-\left\lceil\frac n2\right\rceil
 \qquad(n\geq1),
\]
so $f_c(n)/n\to c-1/2$.
\end{enumerate}
\end{theorem}

We use the following theorem of Alon \cite[Theorem~1.1]{Alon}: for every integer $k\geq2$ and every $\varepsilon>0$, there is $N_0$ such that, whenever $N\geq N_0$ and $X\subseteq\Z/N\Z$ satisfies
\[
 \card X>\left(\frac1k+\varepsilon\right)N,
\]
there is a nonempty $Y\subseteq X$ with $\card Y\leq k$ and $\sum_{y\in Y}y=0$ in $\Z/N\Z$.

The next consequence is the only additive input needed below.

\begin{proposition}\label{prop:bound}
Fix $\rho\geq2$ and $\delta>0$.  For all sufficiently large $E$, the following holds uniformly for every even integer $\tau$ with
\[
 2E\leq\tau\leq\rho E.
\]
If $A\subseteq[1,E]$ and $\tau\notin\Sigma(A)$, then, with
\[
 K=\left\lfloor\frac{3\tau}{2E}\right\rfloor,
\]
we have
\[
 \card A\leq \frac{\tau}{2K}+\delta E.
\]
\end{proposition}

\begin{proof}
Put $q=\tau/2$.  Then $E\leq q\leq\rho E/2$, while $3\leq K\leq3\rho/2$; hence $K$ ranges over a fixed finite set.  We may therefore use Alon's theorem uniformly in $K$, with $\varepsilon=\delta/\rho$, and take $E$ so large that also $K<\delta E/2$.

Suppose that $\card A>q/K+\delta E$.  Reduction modulo $q$ is injective on $[1,E]$, and
\[
 \frac{q}{K}+\frac{\delta q}{\rho}
 \leq \frac{q}{K}+\frac{\delta E}{2}<\card A.
\]
Alon's theorem gives a nonempty $B_1\subseteq A$ with $\card{B_1}\leq K$ and $q\mid\sum_{b\in B_1}b$.  Moreover
\[
 0<\sum_{b\in B_1}b\leq \card{B_1}E\leq KE\leq3q.
\]
The last inequality is strict.  Indeed, equality throughout would force $\card{B_1}=K\geq3$ and every element of $B_1$ to equal $E$, which is impossible.  Thus
\[
 \sum_{b\in B_1}b\in\{q,2q\}.
\]
If this sum is $2q=\tau$, we are done.  Otherwise it is $q$.  After deleting $B_1$ we still have
\[
 \card{A\setminus B_1}>\frac{q}{K}+\delta E-K
 >\frac qK+\frac{\delta E}{2}
 \geq\left(\frac1K+\frac\delta\rho\right)q.
\]
A second application of Alon's theorem gives a nonempty $B_2\subseteq A\setminus B_1$ whose ordinary sum is again $q$ or $2q$.  In the latter case $B_2$ sums to $\tau$; in the former case $B_1\cup B_2$ does.  Hence every set larger than the asserted bound contains a subset summing to $\tau$.
\end{proof}

\begin{proof}[Proof of Theorem~\ref{thm:main}]
First suppose $0<c<1$, and assume for contradiction that
\[
 \frac{f_c(n)}n\longrightarrow L.
\]
For odd $n$, the even elements of $[1,\lfloor cn\rfloor]$ have no subset sum equal to $n$.  Their cardinality is $\lfloor\lfloor cn\rfloor/2\rfloor$, so passage to the odd subsequence gives
\begin{equation}\label{eq:odd}
 L\geq\frac c2.
\end{equation}

Suppose first that $c\leq1/2$, and put
\[
 J=\left\lfloor\frac{3}{2c}\right\rfloor.
\]
Since $\lfloor x\rfloor>x-1$ and $c\leq1/2$, we have $J>1/c$.  Choose $\delta>0$ so that
\[
 \beta:=\frac1{2J}+\delta c<\frac c2.
\]
For sufficiently large even $n$, set $M=\lfloor cn\rfloor$.  Then $2M\leq n\leq(2/c)M$, and Proposition~\ref{prop:bound}, applied with $E=M$ and $\tau=n$, gives
\[
 f_c(n)\leq \frac{n}{2K}+\delta M,
 \qquad
 K=\left\lfloor\frac{3n}{2M}\right\rfloor\geq J.
\]
Consequently $f_c(n)/n\leq\beta$ for all sufficiently large even $n$, and therefore $L\leq\beta<c/2$, contradicting \eqref{eq:odd}.

Now suppose $1/2<c<1$.  Write $d=1-c$ and
\[
 J=\left\lfloor\frac{3}{2d}\right\rfloor.
\]
Since $d<1/2$, the same calculation gives $J>1/d$.  Set
\[
 g=\frac d2-\frac1{2J}>0,
 \qquad
 \delta=\frac{g}{3d},
 \qquad
 \eta=\frac g3.
\]
Then
\begin{equation}\label{eq:gap}
 c-\frac12+\frac1{2J}+\delta d+\eta
 =\frac c2-\frac g3<\frac c2.
\end{equation}
For a sufficiently large even $n$, put $M=\lfloor cn\rfloor$ and $D=n-M$.  If $A\subseteq[1,M]$ avoids $n$, split
\[
 A_0=A\cap[1,D-1],
 \qquad
 A_1=A\cap[D,M].
\]
The interval $[D,M]$ is invariant under $x\mapsto n-x$.  Since $A$ cannot contain both members of a two-element orbit,
\[
 \card{A_1}\leq\left\lfloor\frac{M-D+2}{2}\right\rfloor
 \leq\left(c-\frac12\right)n+1.
\]
Since $dn\leq D<dn+1$, for all sufficiently large $n$ we have $D\leq M$,
$2(D-1)\leq n\leq(2/d)(D-1)$, and $D-1\leq dn$.  Apply Proposition~\ref{prop:bound} to $A_0$ with $E=D-1$, $\tau=n$, and $\rho=2/d$; it gives
\[
 \card{A_0}\leq\frac{n}{2K}+\delta(D-1),
 \qquad
 K=\left\lfloor\frac{3n}{2(D-1)}\right\rfloor\geq J.
\]
Since $D-1\leq dn$, we obtain
\[
 \frac{\card A}{n}
 \leq c-\frac12+\frac1{2J}+\delta d+\frac1n.
\]
Taking $n$ large enough that $1/n\leq\eta$, \eqref{eq:gap} again gives an eventual upper bound strictly below $c/2$ on the even subsequence, contradicting \eqref{eq:odd}.  This proves part~(1).

For part~(2), let $1\leq n\leq M$.  An avoiding set cannot contain $n$.  On $[1,n-1]$, the map
\[
 x\longmapsto\min\{x,n-x\}
\]
is injective on every avoiding set, since two distinct elements with the same image sum to $n$.  Hence an avoiding subset of $[1,M]$ has at most
\[
 \left\lfloor\frac n2\right\rfloor+(M-n)
 =M-\left\lceil\frac n2\right\rceil
\]
elements.  Equality is attained by
\[
 \left[\left\lceil\frac n2\right\rceil,n-1\right]\cup[n+1,M],
\]
which plainly has no subset sum equal to $n$.  Thus
\[
 F(M,n)=M-\left\lceil\frac n2\right\rceil.
\]
Taking $M=\lfloor cn\rfloor$ proves part~(2).
\end{proof}

\begin{thebibliography}{9}

\bibitem{Alon}
N.~Alon,
\emph{Subset sums},
J. Number Theory \textbf{27} (1987), 196--205.

\bibitem{EG}
P.~Erd\H{o}s and R.~L. Graham,
\emph{Old and New Problems and Results in Combinatorial Number Theory},
Monographies de L'Enseignement Math\'ematique, no.~28, Geneva, 1980.

\end{thebibliography}

\end{document}
